%% File: appendices/appendixA.tex
%% Author: Y.U.P. (yashparitkar)
%%
%% Description: Appendix A: exercising the whole flow with no hardware attached.

\documentclass[../manual.tex]{subfiles}

\begin{document}

\chapter{Simulations}
\label{app:nohw}

Everything in this manual except programming a board can be exercised with nothing plugged in. The test suite is how, and it is also the quickest way to find out whether a change you made broke something, since the host side of it runs in milliseconds.

\section{The test suite}
\label{sec:nohw_tests}

From the root of the repository:

\begin{lstlisting}
make sim           # everything
make sim-python    # the host side tests, python3 and nothing else
make sim-c         # the baremetal driver, a C compiler and nothing else
make sim-hdl       # the GHDL simulations
make configure     # just bring the generated core up to date
\end{lstlisting}

\texttt{tests/} is split by what a test \emph{needs} to run rather than by what it covers, which is what makes the three targets independently useful. Table~\ref{tab:nohw_groups} is the whole dependency story.

\begin{table}[H]
    \centering
    \begin{tabularx}{\textwidth}{|l|l|X|}
        \hline
        \rowcolor{gray!40} \textbf{Group} & \textbf{Needs} & \textbf{Covers} \\
        \hline
        \texttt{tests/python/} & python3 & The python driver, every layer of it, against a byte level model of the wrapper \\
        \hline
        \texttt{tests/c/} & A C compiler & The baremetal driver, compiled and run on the host against a stubbed HAL \\
        \hline
        \texttt{tests/hdl/} & GHDL, and the generated core & The core itself, simulated as the generator actually emits it \\
        \hline
    \end{tabularx}
    \caption{The three test groups and what each one needs.}
    \label{tab:nohw_groups}
\end{table}

\texttt{make sim} runs the host side first, and deliberately: those tests take milliseconds, so a broken packet format or register map fails there rather than after the simulations have run. A directory counts as a test if it carries its own \texttt{Makefile}, which keeps the numbering a convention rather than something the build depends on.

Two notes on what is not needed. The python group needs no serial port and no board, because it drives the driver against a model of the wrapper rather than against hardware; \texttt{tests/python/05\_gui} additionally needs PySide6 and skips itself where the widget packages are absent, so \texttt{make sim-python} really does need python3 and nothing else. The C group needs no vendor toolchain, because \texttt{tests/c/00\_reg\_map/stub/} supplies what the PolarFire SoC HAL would have; that stub is also the worked example of the port described in Section~\ref{sec:c_porting}.

\section{What each test proves}
\label{sec:nohw_coverage}

The HDL simulations are the specification of the core's behaviour, one directory per property.

\begin{table}[H]
    \centering
    \begin{tabularx}{\textwidth}{|l|X|}
        \hline
        \rowcolor{gray!40} \textbf{Directory} & \textbf{Proves} \\
        \hline
        \texttt{00\_sim\_ext\_trig} & The external trigger build triggers off \texttt{i\_ext\_trig} \\
        \hline
        \texttt{01\_sim\_axil\_trig} & The internal trigger fires on a condition, here TLAST high \\
        \hline
        \texttt{02\_sim\_axil\_trig\_data} & A masked data pattern triggers on the data it names and not on the rest \\
        \hline
        \texttt{03\_sim\_axil\_cdc} & The two clock domains behave when they are not the same clock \\
        \hline
        \texttt{04\_sim\_axil\_force\_trig} & A second write to ARM\_FT forces the trigger \\
        \hline
        \texttt{06\_sim\_libre\_ila\_uart} & The UART wrapper end to end: configure, arm, trigger and read back over the serial link \\
        \hline
        \texttt{07\_sim\_edge\_test} & Level against edge triggering, TRIG\_CFG bits 1 and 2 \\
        \hline
        \texttt{08\_sim\_addr\_bound\_check} & Address decoding at and past the end of the register map \\
        \hline
        \texttt{09\_sim\_axil\_disarm} & The DISARM register returns an armed core to IDLE \\
        \hline
    \end{tabularx}
    \caption{The HDL simulations.}
    \label{tab:nohw_hdl}
\end{table}

The host side tests cover the two drivers layer by layer.

\begin{table}[H]
    \centering
    \begin{tabularx}{\textwidth}{|l|X|}
        \hline
        \rowcolor{gray!40} \textbf{Directory} & \textbf{Proves} \\
        \hline
        \texttt{python/00\_pkt\_format} & The register reads and writes against a byte level model of the wrapper's packet parser \\
        \hline
        \texttt{python/01\_vcd} & The portmap parsing that turns \texttt{portmap.csv} into a bit layout, and the VCD that comes out of it \\
        \hline
        \texttt{python/02\_cli} & The command line front end end to end, every verb, against that same model \\
        \hline
        \texttt{python/03\_session} & The session layer both front ends drive \\
        \hline
        \texttt{python/04\_trigger} & The translation between per-signal patterns and the flat condition and mask pair \\
        \hline
        \texttt{python/05\_gui} & The graphical front end, driven offscreen against two fake cores \\
        \hline
        \texttt{c/00\_reg\_map} & The baremetal driver against the register map, rather than against itself \\
        \hline
    \end{tabularx}
    \caption{The host side tests.}
    \label{tab:nohw_host}
\end{table}

The HDL group reads \texttt{codegen/gen\_axis/} rather than \texttt{hdl/}, so what is simulated is what the generator actually emits rather than the templates it emits from. Each test Makefile regenerates that directory when the generator, anything in \texttt{hdl/} or anything in \texttt{codegen/templates/} has changed under it, so running one test directly stays as up to date as going through the root Makefile. Section~\ref{sec:codegen_running} covers the one case that stamp cannot see.

\section{Running one on its own}
\label{sec:nohw_single}

Every test directory carries its own Makefile and can be run from where it is:

\begin{lstlisting}
cd tests/hdl/07_sim_edge_test
make sim       # run it
make wave      # run it and open the result in gtkwave
make clean
\end{lstlisting}

\texttt{make wave} is the useful one when a simulation fails: the waveform it opens is the core's own internals, not a capture, so it shows what the trigger comparator and the capture state machine were doing at the moment the test gave up.

The graphical front end has one more entry point, which is not a test at all:

\begin{lstlisting}
cd tests/python/05_gui
make gui       # open the real GUI against two fake cores
make shots     # the same, rendered offscreen to shots/*.png
\end{lstlisting}

\texttt{make gui} starts the real window against a throwaway device store holding two fake cores, so the whole of Chapter~\ref{ch:gui} can be followed without a board attached. It is the fastest way to see what the tool does before committing to building a core into a design.

\end{document}
